% ===========================================================================
\chapter{A database of facts}\label{ch:database}

The engine needs somewhere to keep facts: the ones a program states, and
the far more numerous ones its rules derive. It needs to do three things
with them quickly. Ask whether a fact is already known, so that it is
stored once. Add a new one. And find the facts that might match a pattern,
which is what every rule does all the time.

\section{Relations}

Facts are grouped by predicate. All the \code{parent} facts form one
\emph{relation}, all the \code{ancestor} facts another, as the rows of one
table in a relational database. A relation keeps its facts in the order
they arrived, and the text of each in a set, so that asking whether a fact
is here does not mean looking at every one.

\begin{newfile}{database.nex}
-- database.nex: the facts, kept one relation per predicate.

intern Bindings

-- The facts of one predicate, each once.
class Relation
create
  make(p: String) do
    predicate := p
    facts := []
    keys := #{}
  end
feature
  predicate: String
  facts: Array[Atom]
  -- The text of every fact, for asking quickly whether one is here.
  keys: Set[String]

  count: Integer do
    result := facts.length
  end

  contains(a: Atom): Boolean do
    result := keys.contains(a.text)
  end
\end{newfile}

Adding a fact that is already here does nothing, and \code{add} reports
whether the fact was new. Its postcondition \code{counted} says exactly how
the count changes: by one for a new fact, not at all for a known one.

\begin{continued}{database.nex}
  -- Add `a`, unless it is already here. Is it new?
  add(a: Atom): Boolean
  require
    this_predicate: a.predicate = predicate
    ground: a.is_ground
  do
    result := not contains(a)
    if result then
      facts.add(a)
      keys.add(a.text)
    end
  ensure
    present: contains(a)
    counted: count = old count + (when result then 1 else 0 end)
  end
\end{continued}

\section{Candidates}

Matching a pattern needs a list of facts to try. For now a relation simply
offers all of its facts. In Chapter~\ref{ch:indexing} we will find a
cleverer way that offers fewer, and \code{candidates} is the one place that
will change. Its postcondition says what any way of choosing candidates
must guarantee: \code{complete}, that every fact the pattern matches is
among them. Offering extra facts only wastes a little time, but leaving one
out would make the engine miss answers, silently.

\begin{continued}{database.nex}
  -- The facts that might match `pattern` once `b` is applied to it.
  -- Here that is simply all of them.
  candidates(pattern: Atom, b: Bindings): Array[Atom] do
    result := facts
  ensure
    complete: covers(result, pattern, b)
  end

  -- Is every fact that matches `pattern` under `b` among `found`?
  covers(found: Array[Atom], pattern: Atom, b: Bindings): Boolean do
    result := true
    across facts as f do
      if match_atom(pattern, f, b) /= nil and not found.contains(f) then
        result := false
      end
    end
  end
invariant
  one_key_per_fact: keys.size = facts.length
end
\end{continued}

\section{The database}

The database is a map from predicate to relation. It creates a relation the
first time it sees a predicate, through \code{new_relation}. Making that a
routine of its own looks like a small thing, but it is what will let a
different kind of database, in Chapter~\ref{ch:indexing}, choose a
different kind of relation, without a line of this class changing.

\begin{continued}{database.nex}
-- All the facts known so far.
class Database
create
  make do
    relations := {}
  end
feature
  relations: Map[String, Relation]

  -- A relation for a predicate seen for the first time. A kind of
  -- database that stores its facts differently says so here.
  new_relation(p: String): Relation do
    result := create Relation.make(p)
  ensure
    for_p: result.predicate = p
    empty: result.count = 0
  end

  has(p: String): Boolean do
    result := relations.contains_key(p)
  end

  contains(a: Atom): Boolean do
    result := has(a.predicate) and relations.get(a.predicate).contains(a)
  end

  -- Add a fact. Is it new?
  add(a: Atom): Boolean
  require
    ground: a.is_ground
  do
    if not has(a.predicate) then
      relations.put(a.predicate, new_relation(a.predicate))
    end
    result := relations.get(a.predicate).add(a)
  ensure
    present: contains(a)
    only_grows: count >= old count
  end

  -- Add a fact that `rule` derived with the bindings `b`. Is it new?
  add_derived(a: Atom, rule: Rule, b: Bindings): Boolean
  require
    ground: a.is_ground
    from_rule: substitute(rule.head, b).text = a.text
  do
    result := add(a)
  ensure
    present: contains(a)
  end

  facts_of(p: String): Array[Atom] do
    result := when has(p) then relations.get(p).facts else [] end
  end

  -- The facts that might match `pattern` under `b`.
  candidates(pattern: Atom, b: Bindings): Array[Atom] do
    result := []
    if has(pattern.predicate) then
      result := relations.get(pattern.predicate).candidates(pattern, b)
    end
  end

  count: Integer do
    result := 0
    across relations.values as r do
      result := result + r.count
    end
  end

  -- The predicates with facts, in name order.
  predicates: Array[String] do
    result := relations.keys.sort
  end
end
\end{continued}

\code{add_derived} is the same: today it just adds the fact, but every fact
a rule derives will come through it, with the rule and bindings that
produced it. In Chapter~\ref{ch:explanations} a database that remembers
them will be able to say why each fact holds. Its precondition
\code{from_rule} records what ``derived by this rule with these bindings''
means: substituting the bindings into the rule's head gives the fact.

\section{Trying it out}

\begin{tryit}{try_database.nex}
intern Parser
intern io/Path
intern Database

let db := create Database.make
match parse_program(create Path.make(program.working_directory + "/family.dl").read_text) of
  Ok as ok then
    across ok.value.facts as f do
      db.add(f)
    end
  Err as e then console.print_line(e.error)
end
console.print_line(db.count.to_string + " facts about " + db.predicates.to_string)
match parse_program("parent(alice, bob). parent(bob, zoe).") of
  Ok as ok then
    across ok.value.facts as f do
      console.print_line(f.text + " new? " + db.add(f).to_string)
    end
  Err as e then console.print_line(e.error)
end
console.print_line(db.count.to_string + " facts")
\end{tryit}
\tryoutput{try_database}

The first fact was already known, so adding it again changed nothing.
